% Generated draft: requires mathematical and source review.
\documentclass[11pt]{article}
\usepackage[margin=25mm]{geometry}
\usepackage{fontspec}
\setmainfont{Latin Modern Roman}
\usepackage{amsmath,amssymb,mathtools,unicode-math}
\setmathfont{Latin Modern Math}
\usepackage{enumitem,needspace,xurl,hyperref}
\setlength{\parindent}{0pt}
\setlength{\parskip}{5pt}
\setlength{\emergencystretch}{3em}
\newcommand{\N}{\mathbb{N}}
\newcommand{\Z}{\mathbb{Z}}
\newcommand{\Q}{\mathbb{Q}}
\newcommand{\R}{\mathbb{R}}
\newcommand{\C}{\mathbb{C}}
\newcommand{\F}{\mathbb{F}}
\newcommand{\A}{\mathbb{A}}
\newcommand{\PP}{\mathbb{P}}
\newcommand{\E}{\mathbb{E}}
\newcommand{\eps}{\varepsilon}
\newcommand{\dd}{\,\mathrm{d}}
\begin{document}
\section*{UnsolvedMath ID: 20002358}
\textbf{Source problem code:} catalogue entry
\subsection{Definitions and mathematical statement}
Fix an integral affine Weierstrass model $U=E\setminus\{O\}$ of an elliptic curve over $\Q$. For a set $S$ of rational primes, let $R_S=\Z[S^{-1}]$ and
\[
 U(R_S)=\{P\in U(\Q):x(P),y(P)\in R_S\},\qquad
 H_S=\langle U(R_S)\rangle\le E(\Q).
\]
The source's notation $E(R_S)$ is interpreted as affine coordinate integrality: the Siegel--Mahler finiteness assertion is not about all points on a proper projective model. Write $D(P)$ for the finite set of primes occurring in the denominators of $x(P)$ or $y(P)$, so $P\in U(R_S)$ iff $D(P)\subseteq S$.

\textbf{Question 21, Shlapentokh.} For a rank-two elliptic curve, can an infinite prime set $S$ be chosen with $\operatorname{rank}H_S=1$? If $S$ is infinite of density zero, must its affine integral points remain ``small''? The source cites Siegel--Mahler finiteness for finite $S$, but finite cardinality of the point set is distinct from rank of the subgroup it generates. ``Small'' is not given a precise cardinality or height-growth meaning.

\textbf{Supplied density-zero full-rank construction.} For any fixed rank-$r$ curve/model, choose rational points $P_1,\ldots,P_r$ whose classes form a Mordell--Weil free basis, and let $S_0=\bigcup_iD(P_i)$. If $p_j$ is the $j$th prime, put
\[
 S=S_0\cup\{p_{2^n}:n\ge1\}.
\]
The aid asserts that $S$ is infinite and recursive and has both relative natural prime density and Dirichlet prime density zero, while $P_i\in U(R_S)$ and hence $\operatorname{rank}H_S=r$. Natural density here means $\#\{p\in S:p\le X\}/\pi(X)\to0$, not merely ordinary integer density of a prime subset. Recursiveness is relative to the fixed finite data; no uniform algorithm for finding a Mordell--Weil basis from an arbitrary curve equation is inferred. This construction concerns generated rank and does not prove that $U(R_S)$ itself is infinite or large by a height-counting criterion.

\textbf{Supplied fixed-model obstruction.} On the model labelled 389a1,
\[
 y^2+y=x^3+x^2-2x,
\]
the aid states that $(0,0)$ and $(1,0)$ are independent Mordell--Weil generators. Both are integral and therefore lie in $U(R_S)$ for every $S$. Thus $\operatorname{rank}H_S=2$ for every prime set: no infinite $S$ can give rank one on this model. This refutes a universal assertion for every fixed integral rank-two model; it is not a claim about every alternative model or every refined existence question.

\textbf{Supplied denominator-support criterion.} Let $A=E(\Q)/E(\Q)_{\mathrm{tors}}\cong\Z^2$ and choose a primitive rank-one line $L\subset A$. Then $H_S$ has nonzero free rank contained in $L$ exactly when:
\begin{enumerate}
\item some non-torsion $P$ with class in $L$ has $D(P)\subseteq S$;
\item every $Q$ whose class is outside $L$ satisfies $D(Q)\cap(\mathcal P\setminus S)\ne\varnothing$, where $\mathcal P$ is the set of all primes.
\end{enumerate}
This is a hitting-set problem for denominator supports of points transverse to $L$, together with an allowed nonzero point on $L$. An empty transverse support, supplied by an integral point outside $L$, obstructs it for every $S$. Infinite/density constraints on $S$ are additional requirements not solved by restating this criterion. The source's broad density-zero ``smallness'' question remains distinct from these generated-rank results.
\subsection{Short English statement}
Density zero alone does not force small generated Mordell--Weil rank: the aid gives sparse prime sets containing a finite basis-denominator seed. Its explicit rank-two model already has two integral independent points, so no prime set gives generated rank one there. The exact remaining rank-one selection problem is expressed through denominator supports; other meanings of smallness are unspecified.
\subsection{Sources}
\begin{itemize}
\item Local UnsolvedMath registry, permanent ID 20002358, source code catalogue entry; canonical \texttt{aim-logic-notes.json}, notes[133].
\item AIM 2005 list \emph{Extensions of Hilbert's tenth problem}, Question 21 (Shlapentokh): \url{https://aimath.org/WWN/hilberts10th/hilberts10th.pdf}.
\item Siegel--Mahler theorem named in the original question. The aid supplies the model label 389a1, its equation and generator claim, but no database link or separate proving citation. None is invented.
\item Registry research attempt 1, dated 2026-08-14: partial result, rigor ``proved'', novelty confidence low. The explicitly machine-generated, non-peer-reviewed aid requires independent verification before citation.
\end{itemize}
\end{document}
